\section{How Strongly Is the Claim Supported That Dopamine Is Necessary for the Initiation of Voluntary Action?}\label{02_dopamine_action-how-strongly-is-the-claim-supported-that-dopamine-is-necessary-for-the-initiation-of-voluntary-action}

\begin{quote}
\textbf{Correction note (revised).} An earlier version of this document repeated the popular gloss that dopamine depletion ``preserves hedonic liking'' --- i.e., that pleasure is left intact while only motivation fails. That framing overstates what the evidence shows and silently upgrades Berridge's guarded technical term ``liking'' (an objective, scare-quoted reflex measure he himself distinguishes from conscious pleasure) into a claim about subjective hedonic experience. The human experimental evidence in fact runs the \emph{other} way. This version states the action-initiation findings (which are solid) and the hedonic findings (which do \textbf{not} support ``liking intact'' in humans) at their correct, separate strengths.
\end{quote}

\subsection{TL;DR}\label{02_dopamine_action-tldr}

\begin{itemize}
\tightlist
\item
  \textbf{Solid:} In rodents, near-total absence of dopamine abolishes self-initiated goal-directed action --- animals starve next to available food --- and restoring dopamine specifically to the dorsal striatum (caudate-putamen) rescues it. Dopamine is necessary for the routine, self-initiated, effortful \textbf{execution} of motivated voluntary action under normal conditions. This is established neuroscience.
\item
  \textbf{The ``liking is preserved'' claim does NOT transfer to humans and should not be asserted.} The rodent evidence for preserved ``liking'' is orofacial taste reactivity --- a brainstem reflex preserved even in \emph{decerebrate} rats --- which measures a reflex arc below anything meant by hedonic experience, not felt pleasure. The strongest \emph{human experimental} evidence (D2 blockade RCTs; Parkinson's reward-sensitivity studies) shows dopaminergic disruption \textbf{reduces} reported and physiological hedonic response. So in humans, anhedonia is a real loss of pleasure, not merely a motivational illusion sitting atop intact enjoyment.
\item
  \textbf{Overreach to avoid in both directions:} (a) ``Dopamine \emph{is} pleasure'' --- too strong; liking has a real opioid-mediated component that can be dissociated from dopaminergic wanting under specific manipulations. (b) ``Dopamine is irrelevant to pleasure / liking is intact'' --- also wrong for humans, and it is the version that gets weaponized to dismiss anhedonia as theatrical. The defensible position: \textbf{wanting and liking are partially separable systems that normally operate together; dopamine is the substrate of wanting, effort, and action-initiation, and in humans its disruption also demonstrably degrades felt pleasure.}
\end{itemize}

\begin{center}\rule{0.5\linewidth}{0.5pt}\end{center}

\subsection{Key Findings}\label{02_dopamine_action-key-findings}

\begin{enumerate}
\def\labelenumi{\arabic{enumi}.}
\item
  \textbf{Dopamine-deficient (DD) mice (Zhou \& Palmiter, 1995, \emph{Cell}; Szczypka et al., 1999, \emph{PNAS}) cannot self-initiate feeding and starve.} Targeted inactivation of tyrosine hydroxylase produces mice that are ``severely hypoactive, adipsic, and aphagic'' and die by \textasciitilde3--4 weeks unless rescued by L-DOPA. This is the cleanest demonstration that, in the near-total absence of dopamine, the self-initiated \textbf{execution} of goal-directed action fails.
\item
  \textbf{Viral restoration of dopamine to the caudate-putamen alone rescues feeding and goal-directed action} (Szczypka et al., 2001, \emph{Neuron}; Robinson et al., 2007; Palmiter, 2008, \emph{Ann NY Acad Sci}). Palmiter's interpretation: dorsal-striatal dopamine provides a \emph{permissive signal} that gates the execution of feeding and other goal-directed behaviors.
\item
  \textbf{What the DD mice ``preserve'' is narrower than the popular gloss claims.} They retain a sucrose-over-water preference \emph{once a bout is initiated}, can form a morphine conditioned place preference (Hnasko et al., 2005, \emph{Nature}), and can learn reward locations if locomotion is restored by caffeine (Robinson et al., 2005). These are facts about \textbf{learning} and \textbf{sensory preference}, plus a brainstem \textbf{taste reflex} --- not a measurement of subjective pleasure, and not transferable to a claim that human enjoyment is intact under dopamine loss (see Finding 5).
\item
  \textbf{The ``intact liking'' measure is a decerebrate-survivable reflex.} The Grill \& Norgren (1978) taste-reactivity test --- rhythmic tongue protrusions for sweet, gapes for bitter --- is mediated by brainstem circuitry and is preserved even in decerebrate rats (Grill \& Norgren 1978, part II; \textbf{PMID: 630410}). When dopamine-depleted animals still show sucrose ``liking'' reactions, what is preserved is a reflex that survives removal of the entire forebrain. Calling its persistence ``pleasure is intact'' is a theoretical move, not a measurement --- a point Berridge concedes by writing ``liking'' in scare quotes and describing it as an objective core process \emph{distinct from conscious pleasure} (Berridge \& Kringelbach, 2015, \emph{Neuron} 86:646).
\item
  \textbf{Human experimental evidence indicates dopaminergic disruption REDUCES hedonic experience.} This is the decisive correction to the rodent gloss:

  \begin{itemize}
  \tightlist
  \item
    \textbf{D2 blockade blunts pleasure in healthy volunteers.} Weber et al.~(2023, \emph{J Psychiatr Res}; PMID 36680855): a single 300 mg dose of amisulpride vs.~placebo, n=85, double-blind RCT, produced significantly lower positivity ratings across \emph{all} stimulus categories (p=.026) and lower electrodermal responses (p=.017). Authors: D2 blockade ``can reduce at-the-moment hedonic responsivity in healthy volunteers,'' and antipsychotics likely contribute to clinical anhedonia via D2 antagonism.
  \item
    \textbf{Antipsychotic-induced anhedonia is clinically well-documented} (PMC4826766): first-generation antipsychotics produce apathy, anhedonia, and loss of drive, with depression in 20--40\% of patients --- patient reports describe emotional blunting and reduced capacity for pleasure, not merely reduced pursuit.
  \item
    \textbf{Parkinson's anhedonia is real and sometimes dopamine-unresponsive.} Costello et al.~(2025, \emph{Brain}; PMC12129732): PD depression is ``characterized by lower acceptance of offers, driven by markedly lower incentivization by reward (reward sensitivity).''
  \item
    \textbf{Primate dopamine loss impairs reward behavior writ large} (Wang et al., PMC4600774): heavy DA loss in rhesus impairs morphine CPP --- the opposite of what a clean wanting/liking dissociation predicts.
  \end{itemize}
\item
  \textbf{The ``go-to'' human dissociation studies are weak.} Brauer \& De Wit (1997), Leyton et al.~(2005), Wachtel et al.~(2002) used acute phenylalanine/tyrosine depletion, which produces only modest dopamine reduction, and relied on subjective ``high'' ratings with face-validity problems. Effect sizes are small and replication is inconsistent. The dissociation is cleanest in rodent operant assays; its translation to human subjective experience remains an open research question, not a settled fact (Robinson \& Berridge, 2025, \emph{Annu Rev Psychol} 76:29).
\item
  \textbf{Akinetic mutism (AM) is the closest human analog and is classically a motor-initiation disorder with preserved consciousness.} The Aspen Workgroup (Giacino et al., 2002) considers AM ``primarily a defect in motor initiation rather than in consciousness.'' AM is associated with lesions of mesocorticolimbic dopamine pathways, anterior cingulate, and dorsomedial thalamus, and improves with L-DOPA, bromocriptine, methylphenidate, and amantadine.
\item
  \textbf{Direct evidence that AM patients can be conscious despite total unresponsiveness:} Comanducci et al.~(2023, \emph{Eur J Neurosci} 59:860--873) reported an AM patient behaviorally indistinguishable from vegetative state whose TMS-EEG perturbational complexity index (PCImax = 0.34) exceeded the consciousness threshold (PCI* = 0.31, Casarotto et al., 2016), with fMRI command-following; the patient recovered overt behavior on amantadine.
\item
  \textbf{First-person testimony converges on preserved inner experience under dopaminergic akinesia.} Sacks's \emph{Awakenings} (1973) post-encephalitic patients reported awareness throughout decades of freezing --- one, via spelling board: ``I have no exit. I'm trapped in myself. This stupid body is a prison with windows but no doors.'' Bonifati (2019, \emph{Lancet Neurol}): ``these patients were aware and observant, at varying degrees.'' Parkinson's patients in severe OFF describe the same: intact awareness, abolished initiation. (Contrary historical datum: Cairns's original 1941 AM patient was amnestic for the episodes.)
\item
  \textbf{Dopamine encodes multiple parallel signals, not one thing.} Phasic reward-prediction error (Schultz, Dayan \& Montague, 1997); tonic dopamine setting response vigor / opportunity cost (Niv et al., 2007); incentive salience / ``wanting'' (Berridge); and direct movement signals --- Howe \& Dombeck (2016, \emph{Nature} 535:505) found \textasciitilde100 ms dopaminergic bursts ``capable of triggering locomotion,'' and Azcorra et al.~(2023, \emph{Nat Neurosci} 26:1762) identified an Aldh1a1⁺ nigrostriatal subtype that signals movement without reward and is ``especially vulnerable in Parkinson's disease.'' Salamone \& Correa (2012, \emph{Neuron} 76:470) reframe accumbens dopamine as governing \textbf{effort and behavioral activation} --- overcoming effort costs --- which fits the lived report of dopamine deficiency (``I cannot summon the substrate to engage'') better than ``I want but don't like.''
\end{enumerate}

\begin{center}\rule{0.5\linewidth}{0.5pt}\end{center}

\subsection{Details}\label{02_dopamine_action-details}

\subsubsection{1. The Dopamine-Deficient Mouse Literature: What Is Established vs.~What Is Interpretation}\label{02_dopamine_action-the-dopamine-deficient-mouse-literature-what-is-established-vs.-what-is-interpretation}

\textbf{Established.} TH-knockout mice (Zhou \& Palmiter, 1995, \emph{Cell} 83:1197) that cannot synthesize dopamine become hypoactive, aphagic, and adipsic, and die at \textasciitilde3--4 weeks unless rescued. Daily L-DOPA restores activity and feeding for 6--9 hr per dose (Szczypka et al., 1999, \emph{PNAS} 96:12138). Viral restoration of dopamine synthesis targeted to the caudate-putamen restores feeding; targeted to nucleus accumbens it restores exploration but not feeding (Szczypka et al., 2001, \emph{Neuron} 30:819). Retrograde Cre-mediated restoration in the SNc rescues feeding for months (Hnasko et al., 2006, \emph{PNAS}).

\textbf{Established: action initiation is dissociable from learning and from sensory preference.} DD mice prefer sucrose to water during the bouts they manage (Cannon \& Palmiter, 2003); form morphine CPP (Hnasko et al., 2005); and can learn reward locations if caffeinated during acquisition (Robinson et al., 2005). Palmiter's framing (2008, \emph{Ann NY Acad Sci} 1129:35): dorsal-striatal dopamine is a \emph{permissive signal} gating goal-directed execution --- not the reward signal, not the learning signal.

\textbf{What this does NOT establish.} It does not establish that subjective pleasure is intact. The preserved measures are (a) learning, (b) sensory preference once a bout begins, and (c) a brainstem taste reflex. None is a readout of felt hedonic experience. The leap from ``these survive'' to ``liking/pleasure is intact'' is the register-jump corrected throughout this document.

\textbf{Contested framing.} Salamone \& Correa (2012) argue accumbens dopamine governs \emph{effort allocation} and \emph{behavioral activation} rather than ``motivation'' broadly --- DA-depleted rats still consume free food and pursue low-effort rewards. Niv et al.~(2007) read tonic dopamine as encoding average reward rate / opportunity cost. Both are compatible with Palmiter's permissive gating but reframe the deficit as one of \emph{effortful, self-initiated, costly} action specifically.

\subsubsection{2. Wanting, Liking, and Whether Human Pleasure Survives Dopamine Loss}\label{02_dopamine_action-wanting-liking-and-whether-human-pleasure-survives-dopamine-loss}

\textbf{The dissociation, stated correctly.} Berridge, Venier \& Robinson (1989) showed 6-OHDA dopamine depletion left orofacial ``liking'' reactions to sucrose intact while abolishing feeding. Mu-opioid stimulation of \textasciitilde1 mm³ ``hedonic hotspots'' in rostrodorsal nucleus accumbens shell and posterior ventral pallidum doubles these ``liking'' reactions (Peciña \& Berridge, 2005; Smith \& Berridge, 2007). Amphetamine and electrical/optogenetic stimulation increase ``wanting'' without increasing ``liking.'' \textbf{This genuinely shows wanting and liking are not identical} --- dopamine and opioid systems are separable, and that is a real and useful finding.

\textbf{What the dissociation does NOT show.} It does not show that dopamine is irrelevant to pleasure, nor that human enjoyment is preserved under dopamine loss. Three reasons, all in the literature:

\begin{itemize}
\tightlist
\item
  \textbf{The ``liking'' floor is a brainstem reflex.} Grill \& Norgren (1978) taste reactivity survives decerebration (part II, PMID 630410). It is a component that may feed into conscious pleasure but is not identical to it --- Berridge says so explicitly (Berridge \& Kringelbach, 2015).
\item
  \textbf{Human experimental manipulation of dopamine reduces hedonic response.} Weber et al.~(2023, RCT, n=85): D2 blockade lowered both positivity ratings and electrodermal responses across all stimulus categories. Antipsychotic-induced anhedonia (PMC4826766). Costello et al.~(2025, \emph{Brain}): PD reward insensitivity, in some forms unresponsive to dopamine treatment. Wang et al.~(PMC4600774): rhesus dopamine loss impairs reward behavior generally. The strongest \emph{experimental} human data point \emph{against} ``intact liking.''
\item
  \textbf{The human dissociation studies are weak.} Acute tyrosine depletion (Leyton, Brauer \& De Wit, Wachtel) produces only modest dopamine reduction; small effects; inconsistent replication.
\end{itemize}

\textbf{Phenomenological point (reasonable synthesis).} In humans, pleasure includes anticipation, attention, and being-drawn-toward as constitutive parts, not separable add-ons. Stripping anticipatory/activational dopamine does not isolate ``pure liking''; it produces anhedonia --- the grey deadness of depression and stimulant withdrawal. This is exactly the patient's report, and the report is the stronger datum on a question only a verbal subject can answer.

\textbf{The epistemic point worth flagging for the book.} The field's preference for a clean rodent reflex over convergent human testimony on a question \emph{about subjective experience} is itself a methodological double standard: it elevates an ``objective'' proxy precisely because it discounts the witness. This is not (necessarily) a conspiracy; it is the addict-as-unreliable-narrator prior operating inside citation practice. The defensible claim is about \emph{whose report is allowed to count as evidence} --- which is true and sharp --- not that the rodent work is worthless, which would overshoot.

\subsubsection{3. Akinetic Mutism, Advanced Parkinson's, and Inner Experience}\label{02_dopamine_action-akinetic-mutism-advanced-parkinsons-and-inner-experience}

\textbf{Established.} AM = severely reduced spontaneous movement and speech with preserved wakefulness and brainstem reflexes; arises from lesions of frontal-subcortical circuits with dopaminergic pathways centrally implicated (Arnts et al., 2020, \emph{Neurosci Biobehav Rev} 112:270); responsive to dopaminergic agents.

\textbf{Mechanistic distinctions.} Locked-in syndrome (ventral pontine infarct, descending tracts severed, consciousness uncontroversial --- \emph{not} dopaminergic); vegetative state / unresponsive wakefulness (lacks awareness); AM (a motor-initiation defect, behaviorally resembling VS but mechanistically distinct).

\textbf{Reasonable but contested: consciousness preserved during AM.} Strongest objective evidence is Comanducci et al.~(2023): PCImax 0.34 above the 0.31 threshold, plus fMRI command-following, in a patient behaviorally like VS who later recovered on amantadine. Preserved cognitive ERPs reported earlier (Naccache et al., 2004; Kotchoubey et al., 2003).

\textbf{Where it is uncertain.} AM is heterogeneous --- some lesions spare consciousness, others may not. Cairns's original case was amnestic afterward. \emph{Awakenings} are literary case histories. We cannot ethically deplete human dopamine systems to test cleanly. Defensible claim: in well-studied dopaminergic-AM cases with preserved cortical processing and recovery to full personhood, consciousness was very probably preserved during akinesia --- supporting the dissociation between conscious awareness and the capacity for voluntary action.

\textbf{Overreach.} Asserting that \emph{every} AM or end-stage Parkinson patient is fully conscious during akinesia overshoots; late PD involves spreading Lewy pathology and cholinergic loss, making inner experience genuinely unknowable in some cases.

\subsubsection{4. Strongest Objections to ``Dopamine Is Necessary for Action''}\label{02_dopamine_action-strongest-objections-to-dopamine-is-necessary-for-action}

\textbf{Paradoxical kinesia.} Akinetic Parkinson's patients can suddenly move under existential threat --- Bonanni et al.~(2010, \emph{Neurol Sci}), 14 patients during the L'Aquila earthquake; Schlesinger et al.~(2007), missile attacks. Keefe, Salamone et al.~(1989, \emph{Brain Res}): rats with 98\% striatal dopamine depletion still swim out of an ice bath, and this is \emph{not} blocked by haloperidol or SCH-23390 --- i.e., a non-dopaminergic emergency pathway for action exists. Proponents' response: paradoxical kinesia is rare, transient, externally triggered, and not ordinary self-initiated volition.

\textbf{Dopamine-independent reward.} DD mice form morphine CPP (Hnasko et al., 2005); cocaine CPP in DD mice is serotonin-mediated (Hnasko et al., 2007, \emph{J Neurosci} 27:12484). Salamone \& Correa: DA-depleted animals still consume free food and pursue low-cost rewards.

\textbf{Other neuromodulators.} Norepinephrine, acetylcholine, serotonin, orexin, and the Schiff central-thalamic ``mesocircuit'' all contribute to arousal/motivation. Taylor et al.~(2016, \emph{PNAS} 113:12826): optogenetic VTA-dopamine stimulation restored the righting reflex in 6/6 mice under isoflurane, and D1 antagonist pretreatment blocked it in 6/6 --- dopamine activity is \emph{sufficient} to drive arousal transitions, though alternative arousal pathways mean it is not uniquely required.

\textbf{How necessity-proponents respond.} Mainstream position (Palmiter; Salamone; Berke 2018, \emph{Nat Neurosci}): dopamine is necessary not for \emph{all} motivated behavior in a logically airtight sense, but for \textbf{the normal, self-initiated, effortful execution of goal-directed action under everyday conditions.} Backup pathways exist but do not produce sustained, vigorous, repeated self-initiated action --- without dopamine, DD mice and severe-OFF Parkinson's patients do not feed themselves.

\begin{center}\rule{0.5\linewidth}{0.5pt}\end{center}

\subsection{Recommendations}\label{02_dopamine_action-recommendations}

\textbf{State with high confidence:} \textgreater{} ``In rodents, near-total absence of dopamine abolishes self-initiated goal-directed action --- animals starve next to available food --- and restoring dopamine to the caudate-putamen rescues it (Palmiter, 2008; Szczypka et al., 2001). In humans, akinetic mutism and severe Parkinson's OFF states present a closely analogous picture in which voluntary action is abolished while consciousness and perception appear preserved --- demonstrated objectively by TMS-EEG complexity in at least one AM patient (Comanducci et al., 2023) and consistent with testimony from L-DOPA-awakened post-encephalitic patients (Sacks, 1973). Dopamine is necessary for the routine, self-initiated, effortful execution of voluntary action.''

\textbf{State with confidence --- and use to correct the popular meme:} \textgreater{} ``The claim that dopamine loss `preserves liking' does not hold for human pleasure. The rodent `liking' measure is a brainstem taste reflex that survives decerebration, not a readout of felt pleasure (Grill \& Norgren, 1978; Berridge concedes the distinction). Human experimental dopamine blockade \emph{reduces} hedonic responsivity (Weber et al., 2023), antipsychotics induce anhedonia, and Parkinson's reward insensitivity is real and sometimes dopamine-unresponsive (Costello et al., 2025). In humans, anhedonia is a genuine loss of pleasure, not an illusion masking intact enjoyment.''

\textbf{State with moderate confidence:} \textgreater{} ``Wanting and liking are partially separable --- there is a real opioid-mediated hedonic component dissociable from dopaminergic wanting under specific manipulations. But they normally operate together, and dopamine is the substrate of wanting, effort, behavioral activation, and the conversion of value into initiated action.''

\textbf{Avoid / hedge:} \textgreater{} Strict claims that \emph{all} voluntary behavior requires dopamine (paradoxical kinesia, serotonin-mediated CPP, residual swimming under 98\% depletion are documented exceptions). Also avoid the inverse overclaim that ``dopamine isn't pleasure / liking is intact'' as applied to humans --- it is unsupported and is the precise framing used to dismiss anhedonia as theatrical.

\begin{center}\rule{0.5\linewidth}{0.5pt}\end{center}

\subsection{Caveats}\label{02_dopamine_action-caveats}

\begin{itemize}
\tightlist
\item
  \textbf{Species translation cuts against the gloss, not for it.} The direct ``liking preserved'' evidence is rodent; the human evidence (Weber, Costello, Wang) points the other way. Extrapolating rodent intact-liking to humans is the fragile step, and the contrary human data are stronger than the supporting human data.
\item
  \textbf{Terminological landmine.} ``Motivation,'' ``wanting,'' ``liking,'' ``valence,'' ``pleasure,'' ``agency,'' and ``voluntary action'' are not interchangeable. Much apparent disagreement (Berridge vs.~Salamone vs.~Palmiter) is about which term maps to which dopaminergic signal, not about the facts. Specify which is meant.
\item
  \textbf{``Liking'' in scare quotes ≠ pleasure.} Berridge is explicit that his ``liking'' is an objective core process, not conscious pleasure. Any source dropping the scare quotes and asserting ``pleasure is intact'' has committed a register error.
\item
  \textbf{Consciousness-preserved is inferential.} Even Comanducci's PCI is a \emph{correlate} of capacity for consciousness, not a direct readout; Cairns's amnestic case leaves the universal claim open.
\item
  \textbf{Threshold and compensation.} Behavior is largely unaffected until \textgreater80\% dopamine depletion (Palmiter 2008); the claim concerns \emph{near-total} absence, not modest reductions.
\item
  \textbf{Bibliographic note.} The multimodal AM/PCI/fMRI study is Comanducci, Casarotto, Rosanova, \ldots{} Massimini (2023), \emph{Eur J Neurosci} 59(5):860--873. Arnts is lead author of the 2020 \emph{Neurosci Biobehav Rev} AM review and a 2024 \emph{Sci Rep} DBS-zolpidem case.
\end{itemize}
